\section{Análisis teórico}

Sea \(A\) un arreglo no vacío de \(n\) enteros. Una consulta RMQ sobre
\(A[L..R]\), con \(0\leq L\leq R<n\), devuelve el índice de su elemento
mínimo. Denotamos por \(k=R-L+1\) la longitud del rango y resolvemos los
empates con el \textbf{menor índice}. Para ello, cada candidato se representa
mediante el par \((A[i],i)\): el mínimo lexicográfico compara primero el valor
y luego el índice. Las implementaciones reciben intervalos semiabiertos
\([L,R+1)\); convertir los extremos y extraer el índice de la respuesta cuesta
\(\Theta(1)\).

Usamos \(O(\cdot)\) para cotas superiores y
\(\Theta(\cdot)\) para cotas ajustadas. Suponemos un modelo de palabras de
máquina: valores, índices y pares ocupan espacio constante, y acceder, copiar
y comparar un par cuesta \(\Theta(1)\). El espacio de la estructura incluye
sus copias internas, pero excluye el arreglo externo de entrada; devolver un
índice ocupa \(\Theta(1)\). Las expresiones con \(\log n\), en base dos,
se presentan para \(n\geq2\); con un solo elemento las operaciones RMQ son
constantes. El análisis se deriva del código del curso y de las descripciones
de ambas estructuras \cite{edaa-sparse-table,edaa-segment-tree,cp-sparse-table,cp-segment-tree}.

\subsection{Sparse Table}

Sparse Table almacena respuestas para bloques de longitud potencia de dos.
Si \(S_j[i]\) es el mínimo lexicográfico del bloque que comienza en \(i\)
y tiene \(2^j\) elementos, su construcción utiliza
\[
    S_0[i]=(A[i],i),\qquad
    S_j[i]=\min_{\mathrm{lex}}\bigl(S_{j-1}[i],
                              S_{j-1}[i+2^{j-1}]\bigr).
\]
La implementación reserva \(\ell=\lfloor\log_2 n\rfloor+1\) filas
completas de \(n\) elementos, aunque algunas posiciones no correspondan a
bloques válidos \cite{edaa-sparse-table}. Para una consulta, toma
\(b=\lfloor\log_2 k\rfloor\) y combina
\[
    \min_{\mathrm{lex}}\bigl(S_b[L],S_b[R-2^b+1]\bigr).
\]
Los dos bloques cubren el rango y pueden solaparse. La duplicación de su
intersección no cambia el mínimo, por la idempotencia de esta operación
\cite{cp-sparse-table}.

\begin{itemize}
    \setlength{\itemsep}{2pt}
    \setlength{\parskip}{0pt}
    \item \textbf{Preparación (construcción):} \(\Theta(n\log n)\), tanto
    en mejor como en peor caso. Inicializar las filas ya requiere
    \(\Theta(n\ell)\) trabajo. Además, si
    \(m=\lfloor\log_2 n\rfloor\), se realizan
    \(\sum_{j=1}^{m}(n-2^j+1)=m(n+1)-2^{m+1}+2\)
    combinaciones; cada una cuesta tiempo constante.
    \item \textbf{Mejor caso de consulta RMQ:} \(\Theta(1)\).
    \item \textbf{Peor caso de consulta RMQ:} \(\Theta(1)\). El método
    \texttt{query\_idempotent} realiza dos accesos y una combinación.
    Calcula el nivel con \texttt{std::bit\_width}, que permite obtener
    \(\lfloor\log_2 k\rfloor\) sobre una palabra de ancho fijo sin recorrer
    el rango \cite{cp-sparse-table,cpp-bit-width}.
    \item \textbf{Consulta general:} el método distinto \texttt{query}
    descompone \(k\) en potencias de dos y combina bloques disjuntos.
    Procesa tantos bloques como bits uno tiene \(k\): su costo es
    \(\Theta(\operatorname{popcount}(k))\), con mejor caso \(\Theta(1)\)
    y cota \(O(\log(k+1))\). No es la variante utilizada para medir RMQ.
    \item \textbf{Mantenimiento tras un cambio:} la implementación no ofrece
    \texttt{update}. En el escenario dinámico se reconstruye la tabla,
    con costo \(\Theta(n\log n)\) por modificación, independientemente
    de su posición. Es una estrategia de mantenimiento del estudio, no una
    operación de actualización nativa ni una imposibilidad de diseñar
    otras variantes.
    \item \textbf{Memoria:} \(\Theta(n\log n)\) persistente y
    \(\Theta(1)\) auxiliar por consulta. Para
    \(s=\operatorname{sizeof}(T)\), las filas contienen \(n\ell\) objetos
    de tipo \(T\), cuyo tamaño lógico de datos es \(n\ell s\) bytes,
    además de los metadatos de los vectores y del combinador. La expresión
    que crea las filas utiliza también un vector temporal de \(n\)
    elementos: \(\Theta(n)\) espacio temporal, sin cambiar el orden
    \(\Theta(n\log n)\) del pico total de construcción.
\end{itemize}

El tamaño lógico no debe confundirse con los bytes efectivamente reservados:
estos dependen de las capacidades de los vectores y de la biblioteca usada.
Si durante una reconstrucción coexisten la tabla anterior y la nueva, el pico
puede aumentar, pero conserva el mismo orden asintótico.

\subsection{Segment Tree}

Segment Tree mantiene mínimos de grupos de elementos en una jerarquía binaria
\cite{cp-segment-tree}. La versión del curso es iterativa: reserva un vector
de \(2n\) posiciones, coloca las hojas en \([n,2n)\) y calcula los nodos
internos combinando sus dos hijos. No redondea \(n\) a una potencia de dos
\cite{edaa-segment-tree}.

\begin{itemize}
    \setlength{\itemsep}{2pt}
    \setlength{\parskip}{0pt}
    \item \textbf{Preparación (construcción):} \(\Theta(n)\), tanto en mejor
    como en peor caso. Inicializa \(2n\) posiciones, copia \(n\) hojas y
    calcula exactamente \(n-1\) mínimos de hijos. No realiza \(n\)
    actualizaciones independientes, lo que costaría \(O(n\log n)\).
    \item \textbf{Mejor caso de consulta:} \(\Theta(1)\), por ejemplo al
    consultar un único elemento. En ese caso el intervalo se consume en
    una iteración del recorrido ascendente.
    \item \textbf{Peor caso de consulta:} \(\Theta(\log n)\). Los dos
    extremos ascienden dividiéndose por dos y se incorporan como máximo
    dos nodos por nivel. Existen rangos que requieren recorrer
    \(\Theta(\log n)\) niveles. Más precisamente, para longitud \(k\)
    el bucle queda acotado por \(O(\log(k+1))\), porque la cantidad de
    posiciones pendientes se reduce al menos a la mitad en cada paso.
    Por ello, mantener \(k\) fijo no obliga a observar crecimiento
    logarítmico con \(n\), aunque sí puedan variar los accesos a caché.
    \item \textbf{Actualización puntual:} \(\Theta(\log n)\) en mejor y
    peor caso para \(n\geq2\). Después de escribir una hoja, el método
    \texttt{update} recalcula todos sus ancestros hasta la raíz. No termina
    anticipadamente si el nuevo valor coincide con el anterior o si un
    mínimo intermedio no cambia.
    \item \textbf{Memoria:} \(\Theta(n)\) persistente y \(\Theta(1)\)
    auxiliar por consulta y actualización. Los \(2n\) objetos del vector
    representan \(2ns\) bytes lógicos de datos, más los metadatos y
    cualquier capacidad adicional. La construcción usa \(\Theta(1)\)
    espacio temporal fuera de ese almacenamiento, y el pico total es
    \(\Theta(n)\). Las actualizaciones no crean nuevos nodos.
\end{itemize}

Estas cotas se refieren al código iterativo incorporado, no a otra
implementación recursiva que reserve \(4n\) posiciones. En ambas estructuras,
el contenido de los valores no modifica las reservas principales ni la cantidad
de pasos de construcción; el costo de consulta del árbol depende del rango,
no de encontrar antes un valor especialmente pequeño.

\subsection{Comparación de complejidades}

La tabla resume RMQ con pares de tamaño fijo. La consulta excluye la
construcción y el mantenimiento; en Sparse Table corresponde a
\texttt{query\_idempotent}. La columna de mantenimiento compara reconstruir
Sparse Table con actualizar Segment Tree, no dos constructores.

\begingroup
\centering
\small
\begin{tabular}{@{}lccccc@{}}
\toprule
\textbf{Estructura} & \textbf{Construcción} &
\shortstack{\textbf{Mejor}\\\textbf{consulta}} &
\shortstack{\textbf{Peor}\\\textbf{consulta}} &
\textbf{Mantenimiento} & \shortstack{\textbf{Memoria}\\\textbf{persistente}} \\
\midrule
Sparse Table & \(\Theta(n\log n)\) & \(\Theta(1)\) & \(\Theta(1)\) &
\(\Theta(n\log n)\) & \(\Theta(n\log n)\) \\
Segment Tree & \(\Theta(n)\) & \(\Theta(1)\) & \(\Theta(\log n)\) &
\(\Theta(\log n)\) & \(\Theta(n)\) \\
\bottomrule
\end{tabular}
\par
\endgroup
